% !TeX root = ../main-german.tex
% Eigene Befehle, Anpassungen und Umgebungen; nach packages.tex laden.

% 🇩🇪 wird fuer Tabellen benötigt (z.B. >{centering\RBS}p{2.5cm} erzeugt einen zentrierten 2,5cm breiten Absatz in einer Tabelle
\newcommand{\RBS}{\let\\=\tabularnewline}

% 🇺🇸 To avoid issues with Springer's \mathplus. See also http://tex.stackexchange.com/q/212644/9075
\providecommand\mathplus{+}

% 🇺🇸 from hmks makros.tex - \indexify
\newcommand{\toindex}[1]{\index{#1}#1}

% 🇩🇪 Tipp aus "The Comprehensive LaTeX Symbol List"
\newcommand{\dotcup}{\ensuremath{\,\mathaccent\cdot\cup\,}}

% 🇩🇪 Anstatt $|x|$ $\abs{x}$ verwenden. Die Betragsstriche skalieren automatisch, falls "x" etwas größer sein sollte...
\newcommand{\abs}[1]{\left\lvert#1\right\rvert}

% 🇩🇪 Seitengrößen - Gegen Schusterjungen und Hurenkinder...
\newcommand{\largepage}{\enlargethispage{\baselineskip}}
\newcommand{\shortpage}{\enlargethispage{-\baselineskip}}

\newcommand{\initialism}[1]{%
  \textlcc{#1}\xspace%
}
\newcommand{\OMG}{\initialism{OMG}}
\newcommand{\BPEL}{\initialism{BPEL}}
\newcommand{\BPMN}{\initialism{BPMN}}
\newcommand{\UML}{\initialism{UML}}

% Ein "abstract" ist eine "Kurzfassung", keine "Zusammenfassung"
\addto\captionsngerman{%
  \renewcommand\abstractname{Kurzfassung}%
}

\theorempreskipamount 2ex plus1ex minus0.5ex
\theorempostskipamount 2ex plus1ex minus0.5ex
\theoremstyle{break}
\newtheorem{definition}{Definition}[chapter]

\theorembodyfont{\normalfont}
\newtheorem{example}{Beispiel}[chapter]
\newtheorem{remark}{Bemerkung}[chapter]
\newtheorem{lemma}{Lemma}[chapter]
\newtheorem{theorem}{Theorem}[chapter]
\newtheorem{proposition}{Proposition}[chapter]
\newtheorem{hilfssatz}{Hilfssatz}[chapter]


\theoremstyle{nonumberplain}
\theoremheaderfont{\normalfont\itshape}
\theorembodyfont{\normalfont}
\newtheorem{proof}{Beweis}

\newcommand{\dingcheck}{\ding{51}}
\newcommand{\dingcross}{\ding{55}}


\makeatletter
\AtBeginDocument{\toggletrue{blx@useprefix}}
\AtBeginBibliography{\togglefalse{blx@useprefix}}
\makeatother

% Thin spaces between initials
% http://tex.stackexchange.com/a/11083/9075
\renewrobustcmd*{\bibinitdelim}{\,}

% Keep first and last name together in the bibliography
% http://tex.stackexchange.com/a/196192/9075
\renewcommand*\bibnamedelimc{\addnbspace}
\renewcommand*\bibnamedelimd{\addnbspace}

% Replace last "and" by comma in bibliography
% See http://tex.stackexchange.com/a/41532/9075
\AtBeginBibliography{%
  \renewcommand*{\finalnamedelim}{\addcomma\space}%
}

% enable hyperlinked author names when using \citeauthor
% source: http://tex.stackexchange.com/a/75916/9075
\DeclareCiteCommand{\citeauthor}
{
  \boolfalse{citetracker}%
  \boolfalse{pagetracker}%
  \usebibmacro{prenote}
}
{
  \ifciteindex
  {\indexnames{labelname}}
  {}%
  \printtext[bibhyperref]{\printnames{labelname}}
}
{\multicitedelim}
{\usebibmacro{postnote}}

\newcommand{\textcomment}[2]{\colorbox{yellow!60}{#1}\pdfcomment[color={0.234 0.867 0.211},hoffset=-6pt,voffset=10pt,opacity=0.5]{#2}}

% Small PDF comment
% 1st parameter: comment
\newcommand{\sidecomment}[1]{\pdfcomment[color={0.045 0.278 0.643},voffset=4pt,icon=Note]{#1}}
% Disabled variant - for the final PDF
%\newcommand{\sidecomment}[1]{}

\newcommand{\todo}[1]{TODO!\sidecomment{#1}}

% Changes
%
% 1st parameter: review comment
% 2nd parameter: new text
\newcommand{\change}[2]{{\color{red}#2}\pdfcomment[color={0.234 0.867 0.211},voffset=8pt,opacity=0.5]{#1}}
% Disabled variant - for the final PDF
%\newcommand{\change}[2]{#2}

% Define default commands
\makeatletter
\@ifundefined{missingfigure}{\newcommand{\missingfigure}{... missing figure ...}}{}
\@ifundefined{textcomment}{\newcommand{\textcomment}[2]{#1 \todo{#2}}}{}
\@ifundefined{sidecomment}{\newcommand{\sidecomment}[1]{\marginpar{#1}}}{}
\@ifundefined{todo}{\newcommand{\todo}[1]{\sidecomment{#1}}}{}
\@ifundefined{TODO}{\newcommand{\TODO}[1]{\todo{#1}}}{}
\@ifundefined{todofix}{\newcommand{\todofix}[1]{\todo{#1}}}{}
\@ifundefined{change}{\newcommand{\change}[2]{#1 $\rightarrow$ #2}}{}
\makeatother

% Textmarker (Textfarbe rot)
\newcommand{\textmarker}[1]{{\color{red} #1}\xspace}

% Modified (Text blau)
\newcommand{\modified}[1]{{\color{blue!60!black} #1}\xspace}

\floatname{algorithm}{Algorithmus}
\renewcommand{\listalgorithmname}{Algorithmenverzeichnis}

\newcommand{\commentchar}{\ensuremath{/\mkern-4mu/}}
\algrenewcommand{\algorithmiccomment}[1]{\hfill $\commentchar$ #1}

% \Zcref{...} is short for \zcref[S]{...}: capitalized, unabbreviated name at the beginning of a sentence
\NewDocumentCommand{\Zcref}{s O{} m}{\IfBooleanTF{#1}{\zcref*[S,#2]{#3}}{\zcref[S,#2]{#3}}}

%introduce \powerset - hint by http://matheplanet.com/matheplanet/nuke/html/viewtopic.php?topic=136492&post_id=997377
\DeclareFontFamily{U}{MnSymbolC}{}
\DeclareSymbolFont{MnSyC}{U}{MnSymbolC}{m}{n}
\DeclareFontShape{U}{MnSymbolC}{m}{n}{
  <-6>    MnSymbolC5
  <6-7>   MnSymbolC6
  <7-8>   MnSymbolC7
  <8-9>   MnSymbolC8
  <9-10>  MnSymbolC9
  <10-12> MnSymbolC10
  <12->   MnSymbolC12%
}{}
\DeclareMathSymbol{\powerset}{\mathord}{MnSyC}{180}

\addto\captionsngerman{%
  \renewcommand*{\glossaryname}{Bedeutung}%
}

% Enable hyphenation at other places as the dash.
% Example: applicaiton\hydash specific
\makeatletter
\newcommand{\hydash}{\penalty\@M-\hskip\z@skip}
% Definition of "= taken from http://mirror.ctan.org/macros/latex/contrib/babel-contrib/german/ngermanb.dtx
\makeatother

\newcommand*\Descriptionlabel[1]{%
  \raisebox{0pt}[1ex][0pt]{
    \makebox[\labelwidth][1]{
      \parbox[t]{\labelwidth}{
        \hspace{0pt}\textbf{#1:}}}}
}

\newcommand*\Descriptionlabelx[1]{%
  \parbox[t]{\textwidth}{
    \textbf{#1}\\\mbox{}}
}

\newenvironment{Description}{
  \begin{list}{}{
    \let\makelabel\Descriptionlabelx
    \setlength\labelwidth{1em}
    \setlength\leftmargin{\labelwidth+\labelsep}
  }
}
{
  \end{list}
}

% globally change line spacing of lists
% paralist has suspended development since 10 years.
% enumitem has been updated 2011-09-28

\setlist{partopsep=0pt,itemsep=1pt}

%------------------------------------------------------------------------
% fquote Fancy Quotation environment
% supports empty/optional author

% Use \sloppy to make right-margin easier?
% Set picture units to be relative to font size (em)?
% Use begingroup to rest units afterwards?

% provides \isempty test
\definecolor{quotemark}{gray}{0.7}

%fquote environment with author as optional parameter
%usage: \begin{fquote}quote\end{fquote} or \begin{fquote}[Author]quote\end{fquote}
\newenvironment{fquote}[1][]{%
  \newcommand{\fqauthor}{\relax}
  \ifthenelse{\isempty{#1}}
  {}% do nothing
  {\renewcommand{\fqauthor}{\hfill\textsc{--- #1}}}
  \vspace{1em}
  \begin{list}{}{%
    \setlength{\leftmargin}{0.2\textwidth}
    \setlength{\rightmargin}{0.2\textwidth}
  }
  \item[]%
  \begin{picture}(0,0)(0,0)
    \put(-15,-5){\makebox(0,0){%
        \scalebox{4.5}{\textcolor{quotemark}{\bfseries``}}}%
    }
  \end{picture}\em\ignorespaces%
}{%
  \newline%
  \makebox[0pt][l]{\hspace{0.6\textwidth}%
    \begin{picture}(0,0)(0,0)
      \put(15,10){\makebox(0,0){%
          \scalebox{4.5}{\textcolor{quotemark}{\rmfamily\bfseries''}}}%
      }
    \end{picture}}%
  \fqauthor
  \end{list}
}

%German fquote
%  1 parameter for the author's name, may be empty ("{}")
%  guaranteed German quotes (works with lualatex and babel package)
%  usage: \begin{gfquote}{Author}quote\end{gfquote}
\newenvironment{gfquote}[1]{%
  \newcommand{\fqauthor}{\relax}
  \ifthenelse{\isempty{#1}}
  {}% do nothing
  {\renewcommand{\fqauthor}{\hfill\textsc{\textemdash #1}}}
  \vspace{1em}
  \begin{list}{}{%
    \setlength{\leftmargin}{0.2\textwidth}
    \setlength{\rightmargin}{0.2\textwidth}
  }
  \item[]%
  \begin{picture}(0,0)(0,0)
    \put(-15,-5){\makebox(0,0){%
        \scalebox{4.5}{\textcolor{quotemark}{\bfseries \glqq}}}%
    }
  \end{picture}\em\ignorespaces%
}{%
  \newline%
  \makebox[0pt][l]{\hspace{0.6\textwidth}%
    \begin{picture}(0,0)(0,0)
      \put(15,10){\makebox(0,0){%
          \scalebox{4.5}{\textcolor{quotemark}{\rmfamily\bfseries \grqq}}}%
      }
    \end{picture}}%
  \fqauthor
  \end{list}
}

% Benötigt für eine korrekte Darstellung der Hinweise im erzeugten PDF
\newcount\LTGbeginlineexample
\newcount\LTGendlineexample
\newenvironment{ltgexample}%
{\LTGbeginlineexample=\numexpr\inputlineno+1\relax}%
{\LTGendlineexample=\numexpr\inputlineno-1\relax%
  \tcbinputlisting{%
    listing only,
    listing file=\currfilepath,
    colback=green!5!white,
    colframe=green!25,
    coltitle=black!90,
    coltext=black!90,
    left=8mm,
    title=Zugehöriger \LaTeX{}-Quelltext aus \texttt{\currfilepath},
    listing options={
        frame=none,
        language={[LaTeX]TeX},
        escapeinside={},
        firstline=\the\LTGbeginlineexample,
        lastline=\the\LTGendlineexample,
        firstnumber=\the\LTGbeginlineexample,
        basewidth=.5em,
        % lange Zeilen auch ohne Leerzeichen umbrechen (z.B. \includegraphics[...]{...}); sonst ragen sie aus der Box
        breakatwhitespace=false,
        aboveskip=0mm,
        belowskip=0mm,
        numbers=left,
        xleftmargin=0mm,
        numberstyle=\tiny,
        numbersep=8pt%
      }
  }
}%
